\documentclass{proyectoelectrico}
\usepackage[utf8]{inputenc}
\usepackage[T1]{fontenc}
\usepackage{geometry}
\usepackage{graphicx}
\usepackage{float}
\usepackage{hyperref}
\usepackage[english,indonesian]{babel}

\title{Optimasi Pemrosesan Template LaTeX Menggunakan Pendekatan Berbasis AI}
\author{Academic Support Agent}
\date{\today}

\begin{document}
\maketitle

\begin{abstract}

\end{abstract}

\chapter{PENDAHULUAN}


\chapter{TINJAUAN PUSTAKA}


% Original content (commented out):
% % ------------------------------------
% UNIVERSIDAD DE COSTA RICA
% Facultad de Ingeniería
% Escuela de Ingeniería Eléctrica
% IE0499 - Proyecto Eléctrico
%
% PLANTILLA Y GUÍA DEL TRABAJO ESCRITO
% Versión: v1.0 (marzo 2017)
% ------------------------------------

% Tipo de documento
\documentclass[final]{proyectoelectrico}

% A. PAQUETES Y MACROS ESPECIALES -------
% Paquetes y definiciones que no están incluidos 
% en la clase proyectoelectrico.cls o que son 
% propios del proyecto.
\input{configuracion.tex}

% Tipografía
\usepackage{libertine}
\usepackage{libertinust1math}
\usepackage[T1]{fontenc}
\renewcommand*{\ttdefault}{cmtt}
% ---------------------------------------

% B. DATOS ------------------------------
% Todos los nombres incluyen dos apellidos,
% acentos y signos de puntuación apropiados.
% Revisar las recomendaciones sobre el título
% y los nombres de los profesores en la guía.

% Título del proyecto
\titulo{Título completo del Proyecto Eléctrico}

% Autor (nombre y carné)
\autor{Nikola Tesla Pérez}
\carne{B12345}

% Profesor(a) guía
% (Ing.) Nombre Apellido Apellido(, título)
\guia{Ing. Profesor Guía}

% Profesores lectores 
\lectorA{Ing. Profesor Lector A}
\lectorB{Ing. Profesor Lector B}

% Fecha de entrega del trabajo escrito
\mes{5}		% Número del mes
\ano{1968}	% Formato AAAA
% ---------------------------------------

% C. CONTENIDOS -------------------------

%%%%%%%%%%%%%%%%%%%
\begin{document}
%%%%%%%%%%%%%%%%%%%

\frontmatter

% 1. PORTADA
\portada

% 2. HOJA DE APROBACIÓN
\iffinal
\aprobacion
\fi

% 3. RESUMEN (EN ESPAÑOL E INGLÉS)
\input{contenido/resumen.tex}
\input{contenido/abstract.tex}

% 4. RECONOCIMIENTOS
\iffinal
\input{contenido/reconocimientos.tex}
\fi

% 5. TABLAS DE CONTENIDO, FIGURAS Y TABLAS
\tableofcontents
\listoffigures
\listoftables

% 6. NOMENCLATURA
\input{contenido/nomenclatura.tex}

\mainmatter

% 7. CAPÍTULOS
\input{contenido/1_introduccion.tex}
\input{contenido/2_marco_teorico.tex}
\input{contenido/3_desarrollo.tex}
\input{contenido/4_sobre_LaTeX.tex}
\input{contenido/5_conclusiones.tex}

% 8. APÉNDICES
\appendix
\input{apendices/A_ejemplo.tex}
\input{apendices/B_archivo.tex}
\input{apendices/reglamento_proyecto_electrico.tex}

\backmatter

% 9. BILIOGRAFÍA
\bibliographystyle{plain}
\bibliography{bibliografia/bibliografia.bib}

%%%%%%%%%%%%%%%%%%%
\end{document}
%%%%%%%%%%%%%%%%%%%

% ----------------------------------------------
% Plantilla elaborada por Fabián Abarca Calderón
% ----------------------------------------------

\end{document}